% ============================================================
%  HCMUT-DEE Thesis Kit v2.0.3 — Chương 2
%  Copyright (c) 2026 Nguyen Trong Thang & TS. Nguyen Phuc Khai
%  License: MIT
% ============================================================

\chapter{Soạn Thảo Nội Dung \LaTeX{} Cơ Bản}
\label{ch:basics}

\section{Cấu trúc heading}

Trong \LaTeX{}, các cấp tiêu đề được đánh số tự động:

\begin{lstlisting}[style=latexstyle, caption={Các cấp tiêu đề}]
\chapter{Ten chuong}           % Cap 1 (Chuong)
\section{Ten muc}              % Cap 2 (Muc)
\subsection{Ten muc con}       % Cap 3 (Muc con)
\subsubsection{Ten muc nho}    % Cap 4 (Muc nho hon)
\end{lstlisting}

\textbf{Lưu ý:} Template này đánh số đến cấp \texttt{subsubsection} và hiển thị mục lục đến cấp \texttt{subsection}.

\section{Đoạn văn và xuống dòng}

Trong \LaTeX{}, bạn tạo đoạn văn mới bằng cách \textbf{để trống một dòng}:

\begin{lstlisting}[style=latexstyle, caption={Tạo đoạn văn mới}]
Day la doan van thu nhat. Noi dung tiep tuc
tren cung mot dong.

Day la doan van thu hai. LaTeX tu dong
thut dau dong.
\end{lstlisting}

Để \textbf{xuống dòng} mà không tạo đoạn mới, dùng \verb|\\| hoặc \verb|\newline|.

\section{Danh sách}

\subsection{Danh sách không đánh số (itemize)}

\begin{lstlisting}[style=latexstyle]
\begin{itemize}
  \item Muc thu nhat
  \item Muc thu hai
  \item Muc thu ba
\end{itemize}
\end{lstlisting}

Kết quả:
\begin{itemize}
  \item Mục thứ nhất
  \item Mục thứ hai
  \item Mục thứ ba
\end{itemize}

\subsection{Danh sách đánh số (enumerate)}

\begin{lstlisting}[style=latexstyle]
\begin{enumerate}
  \item Buoc 1: Chuan bi
  \item Buoc 2: Thuc hien
  \item Buoc 3: Kiem tra
\end{enumerate}
\end{lstlisting}

Kết quả:
\begin{enumerate}
  \item Bước 1: Chuẩn bị
  \item Bước 2: Thực hiện
  \item Bước 3: Kiểm tra
\end{enumerate}

\section{Định dạng văn bản}

\begin{table}[H]
\centering
\caption{Các lệnh định dạng văn bản thường dùng}
\label{tab:formatting}
\begin{tabular}{lll}
\toprule
\textbf{Lệnh} & \textbf{Kết quả} & \textbf{Ý nghĩa} \\
\midrule
\verb|\textbf{...}| & \textbf{In đậm} & Bold \\
\verb|\textit{...}| & \textit{In nghiêng} & Italic \\
\verb|\underline{...}| & \underline{Gạch chân} & Underline \\
\verb|\texttt{...}| & \texttt{Monospace} & Code font \\
\verb|\textsuperscript{n}| & x\textsuperscript{n} & Chỉ số trên \\
\verb|\textsubscript{n}| & x\textsubscript{n} & Chỉ số dưới \\
\bottomrule
\end{tabular}
\end{table}

\section{Footnote (Ghi chú chân trang)}

\begin{lstlisting}[style=latexstyle]
Day la mot cau co ghi chu\footnote{Noi dung ghi chu
se xuat hien o cuoi trang.}.
\end{lstlisting}

Đây là một câu có ghi chú\footnote{Nội dung ghi chú sẽ xuất hiện ở cuối trang.}.

\section{Tham chiếu chéo (Cross-reference)}
\label{sec:crossref}

\LaTeX{} cho phép đánh dấu và tham chiếu đến bất kỳ đối tượng nào:

\begin{lstlisting}[style=latexstyle, caption={Cách sử dụng label và ref}]
% Dat nhan (label) cho doi tuong
\chapter{Mo dau}
\label{ch:modau}

\section{Van de nghien cuu}
\label{sec:vande}

% Tham chieu den doi tuong
Nhu da trinh bay o Chuong~\ref{ch:modau},
cu the tai Muc~\ref{sec:vande}...
Hinh~\ref{fig:example} minh hoa...
Bang~\ref{tab:formatting} liet ke...
\end{lstlisting}

\textbf{Quy ước đặt tên label} (khuyến nghị):
\begin{itemize}
  \item Chương: \texttt{ch:ten\_chuong}
  \item Mục: \texttt{sec:ten\_muc}
  \item Hình: \texttt{fig:ten\_hinh}
  \item Bảng: \texttt{tab:ten\_bang}
  \item Công thức: \texttt{eq:ten\_ct}
  \item Listing: \texttt{lst:ten\_code}
\end{itemize}

Ví dụ: Bảng~\ref{tab:formatting} ở trên đã được tham chiếu bằng \verb|\ref{tab:formatting}|.

\section{Viết tiếng Việt trong \LaTeX{}}

Template này đã cấu hình sẵn hỗ trợ tiếng Việt:
\begin{itemize}
  \item Package \texttt{vntex} + \texttt{inputenc} (UTF-8).
  \item Font encoding \texttt{T5} cho tiếng Việt.
  \item Bạn chỉ cần \textbf{gõ tiếng Việt trực tiếp} trong file \texttt{.tex}.
\end{itemize}

\textbf{Lưu ý quan trọng:} Lưu file \texttt{.tex} với encoding \textbf{UTF-8} (hầu hết editor hiện đại đều mặc định UTF-8).
